\pdfbookmark[0]{Notices. Kries: Mathematical Geography. 1814.}{anzeigen-kries-1814}
\begin{center}
\rule{.85\textwidth}{.3pt}\par
Göttingische gelehrte Anzeigen.\quad 13 June 1814.\par
\rule{.85\textwidth}{.3pt}
\end{center}
\hspace*{\parindent}\textit{Textbook of Mathematical Geography by} \textsc{Friedrich Kries}, professor at the Gymnasium in Gotha. With seven copperplates. 236 pages in octavo. Leipzig, published by \textsc{G. J. Göschen}.

In writing this textbook on a science that is of interest in so many ways to anyone with an education, the author aimed to find a suitable middle course. On one side are the scant and superficial outlines usually placed at the beginning of textbooks of political geography, which merely list the main results without justifying or explaining them mathematically. On the other side are the larger works, which assume subtler and less widespread knowledge of higher mathematics. In a work of this kind, one expects only an apt selection from what is already known to be presented in an orderly, thorough and clear way. One does not expect new insights that advance the science itself. The author has indeed achieved that aim. In ten sections he discusses the shape of the Earth in general, the mathematical division of the terrestrial globe and its size, its rotation about its axis and the related phenomena, and how to determine the geographical latitude of a place and lay out a meridian line. He discusses the motion of the Earth around the Sun, the division of the celestial and terrestrial globes in relation to that motion and the phenomena it produces on Earth, and how to determine time and geographical longitude. He then treats the spheroidal shape of the Earth, the making of terrestrial globes and maps, and the use of a terrestrial globe to solve problems in mathematical geography. We can only give our general approval to this organization and choice of topics, although in our view a few more subjects that go unmentioned could have been included here and there, while others, such as the various kinds of map projection, might have been treated a little more briefly. We would, for instance, have liked some guidance on calculating the surface area of individual countries,
